\documentclass[../Odyssey_Project.tex]{subfiles}

\begin{document}
\setcounter{section}{2}
\section{Group Actions}

\subsubsection*{Defining group actions}
\begin{itemize}
    \item A (\emph{left}) \emph{action} of $G$ on a set $X$ is a map $G \times X \to X$, the image of $(g,x)$ denoted $gx$, satisfying $1x = x$ for every $x \in X$ and $(g_1 g_2) x = g_1 (g_2 x)$ for all $g_1, g_2 \in G$, $x \in X$.
    \item If $G$ acts on $X$, we say $X$ is a $G$-set; for any $H \leq G$, $X$ is also an $H$-set by restriction.
    \item For $H \leq G$, the coset space $G/H$ is a $G$-set under left multiplication: $(g, xH) \mapsto gxH$.
\end{itemize}

\begin{proposition}
\label{prop:3.1}
There is a natural bijective correspondence between the set of actions of $G$ on a set $X$ and the set of homomorphisms from $G$ to $S_X$.
\end{proposition}

\begin{proof}
    Suppose first that $G$ acts on $X$. For each $g \in G$, define a function
    \[
        \sigma_g \colon X \to X, \qquad \sigma_g(x) = gx.
    \]
    We claim that $\sigma_g$ is a permutation of $X$. Indeed, $\sigma_{g^{-1}}$ is its inverse, since for every $x \in X$,
    \[
        \sigma_{g^{-1}}(\sigma_g(x)) = g^{-1}(gx) = (g^{-1}g)x = 1x = x
    \]
    and similarly $\sigma_g(\sigma_{g^{-1}}(x)) = x$. Hence $\sigma_g \in S_X$. Now define
    \[
        \Phi \colon G \to S_X, \qquad \Phi(g) = \sigma_g.
    \]
    For $g_1,g_2 \in G$ and $x \in X$, the action law gives
    \[
        \Phi(g_1g_2)(x)
        = \sigma_{g_1g_2}(x)
        = (g_1g_2)x
        = g_1(g_2x)
        = \sigma_{g_1}(\sigma_{g_2}(x)).
    \]
    Thus $\Phi(g_1g_2) = \Phi(g_1)\Phi(g_2)$, so $\Phi$ is a homomorphism.
    Conversely, suppose that $\Phi \colon G \to S_X$ is a homomorphism. Define a rule $G \times X \to X$ by
    \[
        gx = \Phi(g)(x).
    \]
    This rule lands in $X$ because each $\Phi(g)$ is a permutation of $X$. Also,
    \[
        1x = \Phi(1)(x) = x,
    \]
    since $\Phi(1)$ is the identity permutation of $X$. For $g_1,g_2 \in G$ and $x \in X$,
    \[
        (g_1g_2)x
        = \Phi(g_1g_2)(x)
        = (\Phi(g_1)\Phi(g_2))(x)
        = \Phi(g_1)(\Phi(g_2)(x))
        = g_1(g_2x).
    \]
    Therefore this rule is an action of $G$ on $X$.
    Finally, these two constructions undo one another. Starting from an action, the homomorphism sends $g$ to the permutation $x \mapsto gx$, so reconstructing the action gives the same rule. Starting from a homomorphism $\Phi$, the action is defined by $gx = \Phi(g)(x)$, so the permutation attached to $g$ is exactly $\Phi(g)$. Hence the correspondence is bijective.
\end{proof}

\subsubsection*{Faithful actions and Cayley's Theorem}
\begin{itemize}
    \item If $H < G$ with $|G:H| = n$, then the action of $G$ on $G/H$ gives a non-trivial homomorphism $G \to S_n$ via Proposition~\ref{prop:3.1}; if $G$ is simple, this homomorphism must be injective.
    \item The action of $G$ on $X$ is \emph{faithful} (or $G$ acts \emph{faithfully}) if the corresponding homomorphism $G \to S_X$ is injective; equivalently, the only $g \in G$ satisfying $gx = x$ for every $x \in X$ is the identity.
    \item If $G$ acts faithfully on $X$, we call $G$ a \emph{permutation group} on $X$, since $G$ is then imbedded isomorphically in $S_X$.
\end{itemize}

\begin{namedtheorem}[Cayley's Theorem]
\label{thm:cayley}
$G$ is isomorphic with a subgroup of $S_G$; in particular, if $G$ is finite with $|G| = n$, then $G$ is isomorphic with a subgroup of $S_n$.
\end{namedtheorem}

\begin{proof}
    Let $G$ act on itself by left multiplication. In other words, for $g,x \in G$, define
    \[
        gx
    \]
    to be the product of $g$ and $x$ in the group $G$. This is an action: the identity satisfies $1x = x$, and associativity in $G$ gives
    \[
        (g_1g_2)x = g_1(g_2x)
    \]
    for all $g_1,g_2,x \in G$.
    By Proposition~\ref{prop:3.1}, this action gives a homomorphism
    \[
        \Phi \colon G \to S_G.
    \]
    Explicitly, $\Phi(g)$ is the permutation of $G$ given by left multiplication by $g$:
    \[
        \Phi(g)(x) = gx.
    \]
    We now show that $\Phi$ is injective. Suppose $g \in \ker\Phi$. Then $\Phi(g)$ is the identity permutation of $G$, so
    \[
        gx = x
    \]
    for every $x \in G$. Taking $x = 1$, we get $g1 = 1$, and hence $g = 1$. Thus $\ker\Phi = 1$, so $\Phi$ is injective.
    Therefore $G \cong \operatorname{im}\Phi$, and $\operatorname{im}\Phi \leq S_G$. Hence $G$ is isomorphic with a subgroup of $S_G$.
    Finally, if $G$ is finite with $|G| = n$, then $S_G$ is isomorphic with $S_n$, since both are symmetric groups on sets with $n$ elements. Thus $G$ is isomorphic with a subgroup of $S_n$.
\end{proof}

\subsubsection*{$G$-set homomorphisms and orbits}
\begin{itemize}
    \item If $X, Y$ are $G$-sets, $\varphi : X \to Y$ is a \emph{$G$-set homomorphism} if it commutes with the actions: $\varphi(gx) = g\varphi(x)$ for all $g \in G, x \in X$. A bijective such $\varphi$ is a \emph{$G$-set isomorphism}; we then write $X \cong Y$.
    \item For each $x \in X$, the \emph{orbit} of $x$ is $Gx = \{gx \mid g \in G\} \subseteq X$, and the \emph{stabilizer} of $x$ is $G_x = \{g \in G \mid gx = x\} \leq G$.
    \item $Gx$ is itself a $G$-set under the induced action; a subset of $X$ is a $G$-set under the induced action iff it is a union of orbits.
\end{itemize}

\begin{lemma}
\label{lem:3.2}
If $X$ is a $G$-set, then $G_{gx} = g G_x g^{-1}$ for any $g \in G$ and $x \in X$.
\end{lemma}

\begin{proof}
    Let $u \in G$. We show that $u \in G_{gx}$ iff $u \in gG_xg^{-1}$.
    By definition, $u \in G_{gx}$ means that $u$ fixes the element $gx$, so
    \[
        u \in G_{gx}
        \Longleftrightarrow u(gx) = gx.
    \]
    Using the action law, this is equivalent to
    \[
        (ug)x = gx.
    \]
    Now act on both sides by $g^{-1}$. Again using the action law, we get
    \[
        (g^{-1}ug)x = x.
    \]
    This means exactly that $g^{-1}ug \in G_x$. Equivalently, $u \in gG_xg^{-1}$.
    Therefore $G_{gx} = gG_xg^{-1}$.
\end{proof}

\subsubsection*{Transitive actions and the orbit decomposition}
\begin{itemize}
    \item The action of $G$ on $X$ is \emph{transitive} (or $G$ acts \emph{transitively} on $X$) if there is some $x \in X$ such that $Gx = X$, equivalently, if for any $x_1, x_2 \in X$ there exists some $g \in G$ such that $gx_1 = x_2$.
    \item A subset of $X$ is a transitive $G$-set under the induced action iff it consists of a single orbit.
    \item For $H \leq G$, the action of $G$ on $G/H$ is transitive: $(yx^{-1}) xH = yH$.
\end{itemize}

\begin{proposition}
\label{prop:3.3}
Any $G$-set has a unique partition consisting of transitive $G$-sets, namely its partition into orbits.
\end{proposition}

\begin{proof}
    We first show that the orbits form a partition of $X$. Every $x \in X$ lies in its own orbit $Gx$, since $x = 1x$. Now let $x,y \in X$, and suppose that $Gx \cap Gy \neq \varnothing$. Then there are $g_1,g_2 \in G$ such that
    \[
        g_1x = g_2y.
    \]
    Hence
    \[
        y = (g_2^{-1}g_1)x,
    \]
    so $y \in Gx$. It follows that $Gy \subseteq Gx$, since for any $u \in G$,
    \[
        uy = u((g_2^{-1}g_1)x) = (ug_2^{-1}g_1)x \in Gx.
    \]
    By symmetry, we also have $Gx \subseteq Gy$. Therefore $Gx = Gy$. Thus any two orbits are either equal or disjoint, so the orbits partition $X$.
    Each orbit is a transitive $G$-set. Indeed, if $y,z \in Gx$, then $y = g_1x$ and $z = g_2x$ for some $g_1,g_2 \in G$, and
    \[
        (g_2g_1^{-1})y = (g_2g_1^{-1})(g_1x) = g_2x = z.
    \]
    Hence any element of the orbit can be moved to any other element of the same orbit.
    Finally, we prove uniqueness. Suppose $X$ has some partition into transitive $G$-sets, and let $P$ be one piece of this partition. Choose some $x \in P$. Since $P$ is itself a $G$-set, it is closed under the action of $G$. Thus
    \[
        Gx \subseteq P.
    \]
    On the other hand, since $P$ is transitive, every $y \in P$ can be reached from $x$. That is, for every $y \in P$, there is some $g \in G$ such that $gx = y$. Hence every $y \in P$ lies in $Gx$, so
    \[
        P \subseteq Gx.
    \]
    Therefore $P = Gx$. Since every transitive piece $P$ is an orbit, the only possible partition into transitive $G$-sets is the partition into orbits.
\end{proof}

\subsubsection*{Classification of transitive $G$-sets}
\begin{itemize}
    \item Since any $G$-set is a union of transitive $G$-sets, classifying $G$-sets reduces to classifying transitive $G$-sets.
    \item Coset spaces are examples of transitive $G$-sets; we now show every transitive $G$-set is isomorphic to one.
\end{itemize}

\begin{proposition}
\label{prop:3.4}
If $X$ is a transitive $G$-set, then $X \cong G/G_x$ as $G$-sets for any $x \in X$.
\end{proposition}

\begin{proof}
    Fix $x \in X$. Define
    \[
        \varphi \colon G/G_x \to X, \qquad \varphi(gG_x) = gx.
    \]
    We first check that $\varphi$ is well-defined. Suppose $gG_x = hG_x$. Then $h^{-1}g \in G_x$, so $(h^{-1}g)x = x$. Acting by $h$ on both sides gives
    \[
        gx = hx.
    \]
    Hence $\varphi(gG_x) = \varphi(hG_x)$.
    We next show that $\varphi$ is injective. Suppose $\varphi(gG_x) = \varphi(hG_x)$. Then $gx = hx$, so acting by $h^{-1}$ on both sides gives
    \[
        (h^{-1}g)x = x.
    \]
    Thus $h^{-1}g \in G_x$, and therefore $gG_x = hG_x$. Hence $\varphi$ is injective.
    Since $X$ is transitive, for every $y \in X$ there exists some $g \in G$ such that $gx = y$. Thus $y = \varphi(gG_x)$, so $\varphi$ is surjective.
    Finally, we check that $\varphi$ is a $G$-set homomorphism. For $u,g \in G$,
    \[
        \varphi(u(gG_x)) = \varphi(ugG_x) = (ug)x = u(gx) = u\varphi(gG_x).
    \]
    Therefore $\varphi$ is a bijective $G$-set homomorphism, and hence $X \cong G/G_x$ as $G$-sets.
\end{proof}

\begin{corollary}[Orbit-Stabilizer]
\label{cor:3.5}
Let $X$ be a $G$-set. Then $Gx \cong G/G_x$ as $G$-sets for any $x \in X$; in particular, if $G$ is finite, then $|Gx| = |G : G_x|$.
\end{corollary}

\begin{proof}
    The orbit $Gx$ is a transitive $G$-set under the induced action. The stabilizer of $x$ for this restricted action is still $G_x$. Hence, by Proposition~\ref{prop:3.4},
    \[
        Gx \cong G/G_x
    \]
    as $G$-sets.
    If $G$ is finite, then the number of left cosets of $G_x$ in $G$ is $|G : G_x|$. Since $Gx$ is in bijection with $G/G_x$, we have
    \[
        |Gx| = |G/G_x| = |G : G_x|.
    \]
\end{proof}

\subsubsection*{When are transitive $G$-sets isomorphic?}
\begin{itemize}
    \item Having decomposed $G$-sets into transitive ones and identified each with a coset space, the remaining question is when two transitive $G$-sets are isomorphic.
\end{itemize}

\begin{lemma}
\label{lem:3.6}
Let $\varphi : X \to Y$ be a homomorphism of $G$-sets, and let $x \in X$. Then $G_x \leq G_{\varphi(x)}$, and if $\varphi$ is an isomorphism, then $G_x = G_{\varphi(x)}$.
\end{lemma}

\begin{proof}
    Let $g \in G_x$. Then $gx = x$. Since $\varphi$ is a $G$-set homomorphism, we have
    \[
        g\varphi(x) = \varphi(gx) = \varphi(x).
    \]
    Thus $g \in G_{\varphi(x)}$, and therefore $G_x \leq G_{\varphi(x)}$.
    Now suppose that $\varphi$ is an isomorphism. Then $\varphi^{-1} \colon Y \to X$ is also a $G$-set homomorphism. Applying the first part to $\varphi^{-1}$ and the element $\varphi(x) \in Y$, we obtain
    \[
        G_{\varphi(x)} \leq G_{\varphi^{-1}(\varphi(x))} = G_x.
    \]
    Together with the inclusion $G_x \leq G_{\varphi(x)}$, this gives $G_x = G_{\varphi(x)}$.
\end{proof}

\begin{proposition}
\label{prop:3.7}
If $H$ and $K$ are subgroups of $G$, then the $G$-sets $G/H$ and $G/K$ are isomorphic iff $H$ and $K$ are conjugate in $G$.
\end{proposition}

\begin{proof}
    The stabilizer of the coset $H \in G/H$ is exactly $H$. Since every coset of $G/H$ has the form $gH$ for some $g \in G$, Lemma~\ref{lem:3.2} shows that the set of stabilizers arising from the $G$-set $G/H$ is precisely the set of conjugates of $H$. Similarly, the set of stabilizers arising from $G/K$ is precisely the set of conjugates of $K$.
    Suppose first that $\varphi \colon G/H \to G/K$ is a $G$-set isomorphism. By Lemma~\ref{lem:3.6},
    \[
        H = G_H = G_{\varphi(H)}.
    \]
    Now $\varphi(H)$ is some coset $gK \in G/K$, and by Lemma~\ref{lem:3.2} its stabilizer is $gKg^{-1}$. Hence
    \[
        H = G_{\varphi(H)} = G_{gK} = gKg^{-1}.
    \]
    Thus $H$ and $K$ are conjugate in $G$.
    Conversely, suppose that $H = gKg^{-1}$ for some $g \in G$. Then $H$ is the stabilizer of the coset $gK \in G/K$. Since $G/K$ is transitive, Proposition~\ref{prop:3.4} gives
    \[
        G/K \cong G/H
    \]
    as $G$-sets. Therefore $G/H$ and $G/K$ are isomorphic as $G$-sets.
\end{proof}

\subsubsection*{Doubly transitive actions and maximality}
\begin{itemize}
    \item Let $X$ be a $G$-set. We say $X$ is \emph{doubly transitive} (and $G$ acts \emph{doubly transitively} on $X$) if whenever $(x_1, x_2)$ and $(y_1, y_2)$ are elements of $X \times X$ with $x_1 \neq x_2$ and $y_1 \neq y_2$, there exists some $g \in G$ such that $gx_1 = y_1$ and $gx_2 = y_2$.
    \item The natural action of $S_n$ on $\{1, \dots, n\}$ for $n \geq 2$ is doubly transitive. Any doubly transitive $G$-set is transitive.
    \item A proper subgroup $H$ of $G$ is \emph{maximal} if there is no proper subgroup of $G$ that properly contains $H$. Any subgroup of prime index is necessarily maximal by Theorem~\ref{thm:1.6}.
\end{itemize}

\begin{proposition}
\label{prop:3.8}
Let $G$ be a group, let $X$ be a doubly transitive $G$-set, and let $x \in X$. Then $G_x$ is a maximal subgroup of $G$.
\end{proposition}

\begin{proof}
    By Proposition~\ref{prop:3.4}, we have $X \cong G/G_x$ as $G$-sets. Thus we may work with the doubly transitive $G$-set $G/G_x$. Suppose, for a contradiction, that $G_x$ is not maximal. Then there exists a subgroup $K$ such that
    \[
        G_x < K < G.
    \]
    Choose $g \in G$ with $g \notin K$, and choose $k \in K$ with $k \notin G_x$. Since $k \notin G_x$, the cosets $G_x$ and $kG_x$ are distinct. Since $g \notin K$ and $G_x \leq K$, the cosets $G_x$ and $gG_x$ are also distinct.
    Because $G/G_x$ is doubly transitive, there exists some $u \in G$ such that
    \[
        uG_x = G_x
        \quad\text{and}\quad
        u(kG_x) = gG_x.
    \]
    The first equality gives $u \in G_x$, and hence $u \in K$. Since $k \in K$, we have $uk \in K$. The second equality says
    \[
        ukG_x = gG_x,
    \]
    so $g^{-1}uk \in G_x \leq K$. Since both $uk$ and $g^{-1}uk$ lie in $K$, it follows that
    \[
        g = uk(g^{-1}uk)^{-1} \in K,
    \]
    contradicting our choice of $g \notin K$. Therefore no subgroup lies strictly between $G_x$ and $G$, and so $G_x$ is maximal.
\end{proof}

\subsubsection*{Blocks and primitivity}
\begin{itemize}
    \item A non-empty subset $B$ of a transitive $G$-set $X$ is a \emph{block} if $B$ and $gB = \{gx \mid x \in B\}$ are either equal or disjoint for every $g \in G$.
    \item $X$ itself is always a block, and any one-element subset of $X$ is always a block; the transitive $G$-set $X$ is \emph{primitive} if these are the only blocks.
\end{itemize}

\begin{proposition}
\label{prop:3.9}
There is a natural bijective correspondence between the set of blocks of a transitive $G$-set $X$ which contain a given element $x$ and the set of subgroups of $G$ which contain $G_x$.
\end{proposition}

\begin{proof}
    Let $B$ be a block of $X$ containing $x$, and define
    \[
        H_B = \{g \in G \mid gx \in B\}.
    \]
    We first show that $H_B$ is a subgroup of $G$. Since $x \in B$, we have $1x = x \in B$, so $1 \in H_B$. Now let $g,g' \in H_B$. Then $gx \in B$ and $g'x \in B$. Since both $x$ and $gx$ lie in $B$, the sets $B$ and $gB$ meet, and therefore $gB = B$ because $B$ is a block. Hence
    \[
        (gg')x = g(g'x) \in gB = B,
    \]
    so $gg' \in H_B$. Similarly, if $g \in H_B$, then $gx \in B$. Since $g^{-1}(gx) = x \in B$, the sets $g^{-1}B$ and $B$ meet, and therefore $g^{-1}B = B$. Thus $g^{-1}x \in B$, so $g^{-1} \in H_B$. Therefore $H_B \leq G$. Also, if $u \in G_x$, then $ux = x \in B$, so $u \in H_B$. Hence $G_x \leq H_B$.\\
    Thus every block $B$ containing $x$ gives a subgroup $H_B$ containing $G_x$. We now show that this assignment is injective. Suppose $B$ and $B'$ are distinct blocks containing $x$. Choose $y \in B'$ with $y \notin B$. Since $X$ is transitive, there exists some $g \in G$ such that $gx = y$. Then $g \in H_{B'}$, but $g \notin H_B$. Hence $H_B \neq H_{B'}$. \\
    Conversely, let $H$ be a subgroup of $G$ containing $G_x$, and define
    \[
        B_H = Hx = \{hx \mid h \in H\}.
    \]
    Clearly $x \in B_H$, so $B_H$ is non-empty. We claim that $B_H$ is a block. If $g \in H$, then $gB_H = B_H$. Now suppose that $gB_H \cap B_H \neq \varnothing$. Then for some $h_1,h_2 \in H$,
    \[
        gh_1x = h_2x.
    \]
    Hence
    \[
        (h_2^{-1}gh_1)x = x,
    \]
    so $h_2^{-1}gh_1 \in G_x \leq H$. Since $h_1,h_2 \in H$, we get
    \[
        g = h_2(h_2^{-1}gh_1)h_1^{-1} \in H.
    \]
    Therefore, whenever $gB_H$ meets $B_H$, we must have $g \in H$, and hence $gB_H = B_H$. If $gB_H$ does not meet $B_H$, then it is disjoint from $B_H$. Thus $B_H$ is a block.\\
    Finally, we check that this construction gives back the subgroup $H$. By definition,
    \[
        H_{B_H} = \{g \in G \mid gx \in B_H\}.
    \]
    If $g \in H$, then $gx \in Hx = B_H$, so $g \in H_{B_H}$. Thus $H \leq H_{B_H}$. Conversely, if $g \in H_{B_H}$, then $gx = hx$ for some $h \in H$. Therefore
    \[
        (h^{-1}g)x = x,
    \]
    so $h^{-1}g \in G_x \leq H$. Since $h \in H$, we get $g \in H$. Hence $H_{B_H} = H$, proving that every subgroup containing $G_x$ arises from a unique block containing $x$.
\end{proof}

\begin{corollary}
\label{cor:3.10}
Let $X$ be a transitive $G$-set. Then $X$ is primitive iff $G_x$ is a maximal subgroup of $G$ for every $x \in X$.
\end{corollary}

\begin{proof}
    Suppose first that $X$ is primitive, and fix $x \in X$. By Proposition~\ref{prop:3.9}, the blocks of $X$ which contain $x$ correspond exactly to the subgroups of $G$ which contain $G_x$. Since $X$ is primitive, the only blocks containing $x$ are $\{x\}$ and $X$. These correspond to the subgroups $G_x$ and $G$, respectively. Hence there is no subgroup $H$ with
    \[
        G_x < H < G.
    \]
    Therefore $G_x$ is maximal in $G$. \\
    Conversely, suppose that $G_x$ is maximal in $G$ for every $x \in X$. Let $B$ be a block of $X$, and choose some $x \in B$. Again by Proposition~\ref{prop:3.9}, the block $B$ corresponds to a subgroup $H_B$ satisfying
    \[
        G_x \leq H_B \leq G.
    \]
    Since $G_x$ is maximal, either $H_B = G_x$ or $H_B = G$. If $H_B = G_x$, then the corresponding block is $\{x\}$; if $H_B = G$, then the corresponding block is $X$. Thus every block of $X$ is either a one-element subset or all of $X$. Hence $X$ is primitive.
\end{proof}

\subsubsection*{Consequences for primitivity}
\begin{itemize}
    \item If $X$ is a transitive $G$-set, all stabilizers are conjugate by Lemma~\ref{lem:3.2}; so if $G_x$ is maximal for some $x$, it is maximal for every $x$, and Corollary~\ref{cor:3.10} can be modified accordingly.
\end{itemize}

\begin{corollary}
\label{cor:3.11}
Any doubly transitive $G$-set is primitive.
\end{corollary}

\begin{proof}
    This follows from Proposition~\ref{prop:3.8} and Corollary~\ref{cor:3.10}. Indeed, if $X$ is doubly transitive, then Proposition~\ref{prop:3.8} shows that $G_x$ is maximal in $G$ for every $x \in X$. Hence, by Corollary~\ref{cor:3.10}, the $G$-set $X$ is primitive.
\end{proof}

\subsubsection*{Counting via group actions}


\begin{proposition}
\label{prop:3.12}
If $G$ is finite and $H, K \leq G$, then
\[ |HK| = \frac{|H||K|}{|H \cap K|}. \]
\end{proposition}

\begin{proof}
    Let $X = G/K$, and consider $X$ as an $H$-set by restricting the usual left multiplication action of $G$ on $G/K$. The orbit of the coset $K \in G/K$ under this action is $\{hK \mid h \in H\}$.
    The union of the cosets in this orbit is exactly
    \[
        \bigcup_{h \in H} hK = HK.
    \]
    Since distinct left cosets of $K$ are disjoint and each has $|K|$ elements, we have
    \[
        |HK| = |\{hK \mid h \in H\}|\,|K|.
    \]
    The stabilizer of $K$ under the action of $H$ is
    \[
        H_K = \{h \in H \mid hK = K\}.
    \]
    But $hK = K$ iff $h \in K$, and we already have $h \in H$. Hence
    \[
        H_K = H \cap K.
    \]
    By Corollary~\ref{cor:3.5}, this orbit has size
    \[
        |\{hK \mid h \in H\}| = |H : H \cap K| = \frac{|H|}{|H \cap K|}.
    \]
    Therefore
    \[
        |HK| = |\{hK \mid h \in H\}|\,|K| = \frac{|H|}{|H \cap K|}|K| = \frac{|H||K|}{|H \cap K|}.
    \]
\end{proof}

\subsubsection*{Centralizers and conjugacy classes}
\begin{itemize}
    \item For $x \in G$, the \emph{centralizer} of $x$ in $G$ is $C_G(x) = \{g \in G \mid gx = xg\} \leq G$.
    \item For $S \subseteq G$, $C_G(S) = \{g \in G \mid gx = xg \text{ for all } x \in S\} = \bigcap_{x \in S} C_G(x)$, called the centralizer of $S$ in $G$.
    \item $Z(G) = C_G(G)$, and $x \in Z(G)$ iff $C_G(x) = G$.
    \item The \emph{conjugacy class} of $x \in G$ is $\{gxg^{-1} \mid g \in G\}$. By Proposition~\ref{prop:1.10}, the conjugacy class of $\rho \in S_n$ consists of all elements of $S_n$ having the same cycle structure as $\rho$.
    \item More explicitly, if $\rho \in S_n$ has $m_i$ cycles of length $i$ for each $i \geq 1$, then the size of its conjugacy class in $S_n$ is
          \[
              \frac{n!}{\prod_{i \geq 1} i^{m_i}m_i!}.
          \]
          To see the denominator, first arrange the $n$ symbols in an ordered list and then cut this list into $m_i$ blocks of length $i$ for each $i$. This gives $n!$ arrangements. However, each $i$-cycle can be written in $i$ cyclic rotations, and these rotations define the same cycle; this gives a factor of $i^{m_i}$. Also, the $m_i$ cycles of the same length may be permuted among themselves without changing the resulting disjoint cycle decomposition; this gives a factor of $m_i!$. Dividing by all these repetitions gives the formula [24, p.~47].
\end{itemize}

\begin{proposition}
\label{prop:3.13}
The conjugacy classes of $G$ form a partition of $G$, and if $G$ is finite then an element $x \in G$ has $|G : C_G(x)|$ conjugates in $G$.
\end{proposition}

\begin{proof}
    Let $G$ act on itself by conjugation: for $g,x \in G$, define
    \[
        g \cdot x = gxg^{-1}.
    \]
    This is a left action, since $1 \cdot x = x$ and
    \[
        g_1 \cdot (g_2 \cdot x)
        = g_1(g_2xg_2^{-1})g_1^{-1}
        = (g_1g_2)x(g_1g_2)^{-1}
        = (g_1g_2)\cdot x.
    \]
    The orbit of $x \in G$ under this action is
    \[
        \{gxg^{-1} \mid g \in G\},
    \]
    which is exactly the conjugacy class of $x$. Hence the first assertion follows from Proposition~\ref{prop:3.3}, since the orbits of any group action form a partition.
    If $G$ is finite, then Corollary~\ref{cor:3.5} says that the size of the orbit of $x$ is the index of the stabilizer of $x$. The stabilizer is
    \[
        G_x = \{g \in G \mid gxg^{-1} = x\}.
    \]
    But $gxg^{-1}=x$ iff $gx=xg$, so
    \[
        G_x = C_G(x).
    \]
    Therefore the conjugacy class of $x$ has size $|G : C_G(x)|$.
\end{proof}

\subsubsection*{Normalizers and conjugates of subgroups}
\begin{itemize}
    \item A subgroup of $G$ is normal iff it is a union of conjugacy classes; such a union is in fact disjoint.
    \item The order of a normal subgroup of a finite group $G$ must be a sum of orders of conjugacy classes of $G$.
    \item For $H \leq G$, the \emph{normalizer} of $H$ in $G$ is $N_G(H) = \{g \in G \mid gHg^{-1} = H\}$. We have $N_G(H) \leq G$ and $H \trianglelefteq N_G(H)$; $N_G(H)$ is the largest subgroup of $G$ in which $H$ is normal, so $N_G(H) = G$ iff $H \trianglelefteq G$.
\end{itemize}

\begin{proposition}
\label{prop:3.14}
A subgroup $H$ of a finite group $G$ has exactly $|G : N_G(H)|$ conjugates in $G$. In particular, the number of conjugates of $H$ in $G$ divides $|G : H|$ and is equal to $1$ iff $H \trianglelefteq G$.
\end{proposition}

\begin{proof}
    Let $\mathcal{P}(G)$ be the set of subsets of $G$, and let $G$ act on $\mathcal{P}(G)$ by conjugation:
    \[
        g \cdot S = gSg^{-1}
    \]
    for $g \in G$ and $S \subseteq G$. This is a left action, since $1S1^{-1} = S$ and
    \[
        g_1 \cdot (g_2 \cdot S)
        = g_1(g_2Sg_2^{-1})g_1^{-1}
        = (g_1g_2)S(g_1g_2)^{-1}
        = (g_1g_2)\cdot S.
    \]
    The orbit of $H$ under this action is
    \[
        \{gHg^{-1} \mid g \in G\},
    \]
    which is precisely the set of conjugates of $H$ in $G$. The stabilizer of $H$ is
    \[
        \{g \in G \mid gHg^{-1} = H\} = N_G(H).
    \]
    Therefore, by Corollary~\ref{cor:3.5}, the number of conjugates of $H$ in $G$ is
    \[
        |G : N_G(H)|.
    \]
    Since $H \leq N_G(H) \leq G$, we have
    \[
        |G : H| = |G : N_G(H)|\,|N_G(H) : H|.
    \]
    Hence the number of conjugates of $H$ divides $|G : H|$.
    Finally, $H$ has exactly one conjugate iff $|G : N_G(H)| = 1$, iff $N_G(H)=G$. This is equivalent to $gHg^{-1}=H$ for every $g \in G$, which is exactly the condition that $H \trianglelefteq G$.
\end{proof}

\subsubsection*{Double cosets}
\begin{itemize}
    \item Let $G$ be a group and let $H, K$ be subgroups of $G$. A \emph{double coset} of $H$ and $K$ in $G$ is a set $HxK = \{hxk \mid h \in H, k \in K\}$ for some $x \in G$.
    \item Any two double cosets are either disjoint or equal, as is the case with ordinary cosets.
\end{itemize}

\begin{proposition}
\label{prop:3.15}
If $G$ is finite and $H, K \leq G$, then for any $x \in G$ we have
\[ |HxK| = \frac{|H||K|}{|H \cap xKx^{-1}|} = \frac{|H||K|}{|x^{-1}Hx \cap K|}. \]
\end{proposition}

\begin{proof}
    Consider $G/K$ as an $H$-set by restricting the usual left multiplication action. The orbit of $xK \in G/K$ under this action is
    \[
        \{hxK \mid h \in H\}.
    \]
    The union in $G$ of the cosets in this orbit is exactly $HxK$. Since distinct left cosets of $K$ are disjoint and each has $|K|$ elements, we have
    \[
        |HxK| = |\{hxK \mid h \in H\}|\,|K|.
    \]
    The stabilizer of $xK$ under the action of $H$ is
    \[
        \{h \in H \mid hxK = xK\}.
    \]
    Now
    \[
        hxK = xK
        \Longleftrightarrow x^{-1}hx \in K
        \Longleftrightarrow h \in xKx^{-1}.
    \]
    Thus the stabilizer is $H \cap xKx^{-1}$. By Corollary~\ref{cor:3.5},
    \[
        |\{hxK \mid h \in H\}| = |H : H \cap xKx^{-1}|
        = \frac{|H|}{|H \cap xKx^{-1}|}.
    \]
    Therefore
    \[
        |HxK|
        = \frac{|H|}{|H \cap xKx^{-1}|}|K|
        = \frac{|H||K|}{|H \cap xKx^{-1}|}.
    \]
    Finally, conjugating the denominator by $x^{-1}$ gives
    \[
        x^{-1}(H \cap xKx^{-1})x = x^{-1}Hx \cap K.
    \]
    Conjugation is a bijection, so these two intersections have the same order. Hence
    \[
        \frac{|H||K|}{|H \cap xKx^{-1}|}
        =
        \frac{|H||K|}{|x^{-1}Hx \cap K|}.
    \]
\end{proof}

\subsection*{Exercises}

\begin{exercise}
\label{ex:3.1}
Show that a finite simple group whose order is at least $r!$ cannot have a proper subgroup of index $r$.
\end{exercise}
\begin{solution}
    Suppose, for a contradiction, that $H < G$ and $|G:H| = r$.
    Let $G$ act on the coset space $G/H$ by left multiplication. Since $G/H$ has $r$ elements, Proposition~\ref{prop:3.1} gives a homomorphism
    \[
        \Phi \colon G \to S_{G/H} \cong S_r.
    \]
    This is the same permutation-representation idea as in \nameref{thm:cayley}, with the coset space $G/H$ in place of $G$ itself.
    Now $\ker\Phi \trianglelefteq G$. Since $G$ is simple, either $\ker\Phi = G$ or $\ker\Phi = 1$.
    If $\ker\Phi = G$, then every $g \in G$ fixes every coset. In particular,
    \[
        gH = H
    \]
    for every $g \in G$. Hence every $g \in G$ lies in $H$, so $G = H$, contradicting that $H$ is proper. Therefore $\ker\Phi \neq G$, and so $\ker\Phi = 1$.
    Thus $\Phi$ is injective, so $G$ is isomorphic to a subgroup of $S_r$. Hence
    \[
        |G| \leq |S_r| = r!.
    \]
    By hypothesis $|G| \geq r!$, so $|G| = r!$. Therefore $\Phi(G)$ is a subgroup of $S_r$ with the same order as $S_r$, and hence $\Phi(G) = S_r$. Thus $G \cong S_r$.
    For $r \geq 3$, this is impossible because $A_r \trianglelefteq S_r$ is a non-trivial proper normal subgroup. Hence no such subgroup $H$ exists.
    As stated, the exercise also needs the condition $r \geq 3$: for $r = 2$, the group $C_2$ is simple, has order $2 = 2!$, and has the proper subgroup $1$ of index $2$.
\end{solution}

\begin{exercise}
\label{ex:3.2}
Show that a group $G$ acts doubly transitively on a set $X$ iff $G_x$ acts transitively on $X - \{x\}$ for every $x \in X$. (Here we must assume that $X$ has more than two elements.)
\end{exercise}
\begin{solution}
    Suppose first that $G$ acts doubly transitively on $X$, and fix $x \in X$. We show that $G_x$ acts transitively on $X - \{x\}$.
    First, $G_x$ preserves $X - \{x\}$. Indeed, if $g \in G_x$ and $y \neq x$, then $gy \neq x$; otherwise $gy = x$, and applying $g^{-1}$ gives $y = g^{-1}x = x$, a contradiction.
    Now let $y,z \in X - \{x\}$. Then $(x,y)$ and $(x,z)$ are ordered pairs of distinct elements of $X$. Since $G$ acts doubly transitively on $X$, there exists $g \in G$ such that
    \[
        gx = x
        \quad\text{and}\quad
        gy = z.
    \]
    The first equality says that $g \in G_x$, and the second says that this element of $G_x$ sends $y$ to $z$. Hence $G_x$ acts transitively on $X - \{x\}$.
    Conversely, suppose that $G_x$ acts transitively on $X - \{x\}$ for every $x \in X$. Let $(x_1,x_2)$ and $(y_1,y_2)$ be ordered pairs of distinct elements of $X$. We must find some $g \in G$ such that $gx_1 = y_1$ and $gx_2 = y_2$.
    We first find an element $a \in G$ such that $ax_1 = y_1$. If $x_1 = y_1$, take $a = 1$. If $x_1 \neq y_1$, choose
    \[
        w \in X - \{x_1,y_1\},
    \]
    which is possible since $X$ has more than two elements. By hypothesis, $G_w$ acts transitively on $X - \{w\}$. Since $x_1,y_1 \in X - \{w\}$, there exists $a \in G_w$ such that
    \[
        ax_1 = y_1.
    \]
    Thus, in either case, we have some $a \in G$ with $ax_1 = y_1$.
    Since $a$ acts as a bijection on $X$ and $x_2 \neq x_1$, we have $ax_2 \neq ax_1 = y_1$. Hence $ax_2 \in X - \{y_1\}$. Also $y_2 \in X - \{y_1\}$, since $y_2 \neq y_1$. By hypothesis, $G_{y_1}$ acts transitively on $X - \{y_1\}$, so there exists $b \in G_{y_1}$ such that
    \[
        b(ax_2) = y_2.
    \]
    Since $b \in G_{y_1}$, we also have $by_1 = y_1$. Now set $g = ba$. Then
    \[
        gx_1 = b(ax_1) = by_1 = y_1
    \]
    and
    \[
        gx_2 = b(ax_2) = y_2.
    \]
    Therefore $G$ acts doubly transitively on $X$.
\end{solution}

\begin{exercise}
\label{ex:3.3}
Show directly from the definitions (that is, without reference to Propositions~\ref{prop:3.8} and~\ref{prop:3.9}) that a doubly transitive $G$-set is primitive.
\end{exercise}
\begin{solution}
    Let $B$ be a block of $X$. We show that $B$ is either a one-element subset of $X$ or all of $X$.
    If $B$ has exactly one element, then $B$ is one of the trivial blocks. Suppose instead that $B$ has at least two elements, and choose distinct elements $x,y \in B$. We will show that $B = X$.
    Let $z \in X$. If $z = x$, then $z \in B$. So suppose that $z \neq x$. Since $x \neq y$ and $x \neq z$, the ordered pairs
    \[
        (x,y)
        \quad\text{and}\quad
        (x,z)
    \]
    are pairs of distinct elements of $X$. Since $G$ acts doubly transitively on $X$, there exists $g \in G$ such that
    \[
        gx = x
        \quad\text{and}\quad
        gy = z.
    \]
    Now $x \in B$ and $gx = x \in B$, so $gB$ meets $B$. Since $B$ is a block, this forces $gB = B$. But $y \in B$, so $gy \in gB$. Therefore
    \[
        z = gy \in gB = B.
    \]
    Hence every $z \in X$ lies in $B$, so $B = X$.
    Thus every block of $X$ is either a one-element subset or all of $X$. Therefore $X$ is primitive.
\end{solution}

\begin{exercise}
\label{ex:3.4}
Let $G$ be the subgroup of $S_5$ generated by $(1\ 2\ 3\ 4\ 5)$, and let $G$ act on $X = \{1,2,3,4,5\}$ in the canonical way. Show that this action is primitive, but not doubly transitive.
\end{exercise}
\begin{solution}
    Let $\sigma = (1\ 2\ 3\ 4\ 5)$, so $G = \langle \sigma \rangle$ has order $5$.
    First we show that the action is primitive. The action is transitive, since powers of $\sigma$ move any element of $X$ to any other element of $X$. Fix $x \in X$. No non-identity power of $\sigma$ fixes $x$, since $\sigma^k$ moves $x$ by $k$ steps around the $5$-cycle for $1 \leq k \leq 4$. Hence
    \[
        G_x = 1.
    \]
    Since $G$ has prime order, its only subgroups are $1$ and $G$. Therefore $G_x = 1$ is maximal in $G$. By Corollary~\ref{cor:3.10}, the action of $G$ on $X$ is primitive.
    We now show that the action is not doubly transitive. If it were doubly transitive, then there would exist some $g \in G$ sending the ordered pair $(1,2)$ to the ordered pair $(1,3)$; that is,
    \[
        g1 = 1
        \quad\text{and}\quad
        g2 = 3.
    \]
    But the only element of $G$ which fixes $1$ is the identity, and the identity sends $2$ to $2$, not to $3$. Thus no such $g$ exists. Therefore the action is not doubly transitive.
\end{solution}

\begin{exercise}
\label{ex:3.5}
Let $N \trianglelefteq G$, and let $y \in N$. Show that the conjugacy class of $y$ in $G$ is a union of conjugacy classes of the group $N$. Show further that there is a bijective correspondence between the conjugacy classes of $N$ which comprise the conjugacy class of $y$ in $G$ and the cosets of $N C_G(y)$ in $G$.
\end{exercise}
\begin{solution}
    Let
    \[
        \mathcal{C} = \{gyg^{-1} \mid g \in G\}
    \]
    be the conjugacy class of $y$ in $G$. Since $N \trianglelefteq G$ and $y \in N$, every element $gyg^{-1}$ lies in $N$. Hence $\mathcal{C} \subseteq N$.
    Now let $z \in \mathcal{C}$. By definition of $\mathcal{C}$, there is some $g \in G$ such that
    \[
        z = gyg^{-1}.
    \]
    If $n \in N$, then
    \[
        nzn^{-1}
        = n(gyg^{-1})n^{-1}
        = (ng)y(ng)^{-1}.
    \]
    Since $ng \in G$, this is again an element of $\mathcal{C}$. Therefore, whenever $\mathcal{C}$ contains one element $z$, it contains every $N$-conjugate of $z$. Thus $\mathcal{C}$ is a union of conjugacy classes of $N$.
    For the second assertion, set $C = C_G(y)$. Since $N \trianglelefteq G$, the product $NC$ is a subgroup of $G$. We define a map from the left cosets of $NC$ in $G$ to the $N$-conjugacy classes contained in $\mathcal{C}$ by
    \[
        \Psi \colon G/NC \to \{\text{$N$-conjugacy classes contained in }\mathcal{C}\},
        \qquad
        \Psi(gNC) = \{n(gyg^{-1})n^{-1} \mid n \in N\}.
    \]
    In words, the coset $gNC$ is sent to the $N$-conjugacy class of $gyg^{-1}$.
    We first check that this is well-defined. Suppose $hNC = gNC$. Then $h = gmc$ for some $m \in N$ and $c \in C$. Since $c \in C_G(y)$, we have $cyc^{-1} = y$. Hence
    \[
        hyh^{-1}
        = gmcyc^{-1}m^{-1}g^{-1}
        = gmym^{-1}g^{-1}.
    \]
    But
    \[
        gmym^{-1}g^{-1}
        =
        (gmg^{-1})(gyg^{-1})(gmg^{-1})^{-1}.
    \]
    Since $N \trianglelefteq G$, we have $gmg^{-1} \in N$. Therefore $hyh^{-1}$ is $N$-conjugate to $gyg^{-1}$. Thus $hNC$ and $gNC$ give the same $N$-conjugacy class, so $\Psi$ is well-defined.
    The map $\Psi$ is onto. Indeed, take any $N$-conjugacy class contained in $\mathcal{C}$. It contains some element $gyg^{-1}$ with $g \in G$, so it is exactly the image of the coset $gNC$.
    Finally, we show that $\Psi$ is one-to-one. Suppose
    \[
        \Psi(gNC) = \Psi(hNC).
    \]
    Then $gyg^{-1}$ and $hyh^{-1}$ are $N$-conjugate, so there exists $n \in N$ such that
    \[
        hyh^{-1}
        = n(gyg^{-1})n^{-1}
        = (ng)y(ng)^{-1}.
    \]
    This means that $h$ and $ng$ conjugate $y$ to the same element. Equivalently,
    \[
        ((ng)^{-1}h)y((ng)^{-1}h)^{-1} = y.
    \]
    Hence $(ng)^{-1}h \in C_G(y) = C$. So $h = ngc$ for some $c \in C$. Since $N$ is normal in $G$, we have $g^{-1}ng \in N$, and therefore
    \[
        h = ngc = g(g^{-1}ng)c \in gNC.
    \]
    Thus $hNC = gNC$. Therefore $\Psi$ is injective.
    Hence the $N$-conjugacy classes which make up the $G$-conjugacy class of $y$ are in bijective correspondence with the cosets of $NC_G(y)$ in $G$.
\end{solution}

\begin{exercise}
\label{ex:3.6}
Let $G$ be a finite group, and let $r$ be the number of conjugacy classes of $G$. Show that $|\{(a,b) \in G \times G \mid ab = ba\}| = r|G|$.
\end{exercise}
\begin{solution}
    We count the same set in a slower way. Fix $a \in G$. The elements $b \in G$ such that $ab = ba$ are exactly the elements of the centralizer $C_G(a)$. Hence the number of commuting pairs is
    \[
        \sum_{a \in G} |C_G(a)|.
    \]
    We now group this sum by conjugacy classes. Let $K$ be one conjugacy class of $G$, and choose a representative $x \in K$.
    If $a \in K$, then $a = gxg^{-1}$ for some $g \in G$. Conjugation by $g$ gives a bijection
    \[
        C_G(x) \to C_G(a), \qquad u \mapsto gug^{-1},
    \]
    so $|C_G(a)| = |C_G(x)|$. Thus, as $a$ runs through the elements of $K$, the number $|C_G(a)|$ is always the same number $|C_G(x)|$.
    Therefore the part of the sum coming from this one conjugacy class is
    \[
        \sum_{a \in K} |C_G(a)|
        =
        \underbrace{|C_G(x)| + \cdots + |C_G(x)|}_{|K|\text{ times}}
        =
        |K|\,|C_G(x)|.
    \]
    By Proposition~\ref{prop:3.13},
    \[
        |K| = |G : C_G(x)|.
    \]
    Since $G$ is finite, this means
    \[
        |K|\,|C_G(x)|
        =
        |G : C_G(x)|\,|C_G(x)|
        =
        |G|.
    \]
    Hence each conjugacy class contributes exactly $|G|$ to the total sum. Since there are $r$ conjugacy classes, we get
    \[
        |\{(a,b) \in G \times G \mid ab = ba\}|
        =
        r|G|.
    \]
\end{solution}

\begin{exercise}
\label{ex:3.7}
Show that $C_G(gxg^{-1}) = g C_G(x) g^{-1}$ for any elements $g$ and $x$ of a group $G$.
\end{exercise}
\begin{solution}
    Let $G$ act on itself by conjugation, so that $u \cdot x = uxu^{-1}$ for $u,x \in G$.
    Under this action, the stabilizer of $x$ is
    \[
        G_x = \{u \in G \mid uxu^{-1} = x\}.
    \]
    But $uxu^{-1}=x$ iff $ux=xu$, so
    \[
        G_x = C_G(x).
    \]
    Also, the image of $x$ under $g$ is
    \[
        g \cdot x = gxg^{-1}.
    \]
    Therefore Lemma~\ref{lem:3.2}, applied to this conjugation action, gives
    \[
        C_G(gxg^{-1})
        =
        G_{g \cdot x}
        =
        gG_xg^{-1}
        =
        gC_G(x)g^{-1}.
    \]
\end{solution}

\begin{exercise}
\label{ex:3.8}
Let $n \geq 5$, and assume that $A_n$ is simple. (A proof of this fact is outlined in Exercise~\ref{ex:7.8}.) Use Exercise~\ref{ex:3.1} to show that $S_n$ has no proper subgroup of index less than $n$ other than $A_n$.
\end{exercise}
\begin{solution}
    Let $H < S_n$ be a proper subgroup with
    \[
        |S_n : H| < n.
    \]
    We show that $H = A_n$.
    Suppose, for a contradiction, that $H \neq A_n$, and set
    \[
        K = H \cap A_n.
    \]
    We first check that $K$ is a proper subgroup of $A_n$. If $K = A_n$, then $A_n \leq H$. Since $H \neq A_n$, the subgroup $H$ must contain some odd permutation. But then $H$ contains all even permutations and at least one odd permutation. Multiplying that odd permutation by the elements of $A_n$ gives all odd permutations, so $H = S_n$, contradicting that $H$ is proper. Hence
    \[
        K < A_n.
    \]
    We now show that $K$ has index less than $n$ in $A_n$. There are two cases.
    First suppose that $H \leq A_n$. Then $K = H$, and
    \[
        |S_n : H|
        =
        |S_n : A_n|\,|A_n : H|
        =
        2|A_n : K|.
    \]
    Therefore $|A_n : K| < n$.
    Now suppose that $H \nleq A_n$. Then $H$ contains an odd permutation, so
    \[
        HA_n = S_n.
    \]
    By Proposition~\ref{prop:3.12},
    \[
        |S_n|
        =
        |HA_n|
        =
        \frac{|H||A_n|}{|H \cap A_n|}
        =
        \frac{|H||A_n|}{|K|}.
    \]
    Dividing by $|H|$ gives
    \[
        |S_n : H|
        =
        |A_n : K|.
    \]
    Hence again $|A_n : K| < n$.
    Thus, in either case, $K$ is a proper subgroup of $A_n$ of index
    \[
        s = |A_n : K| < n.
    \]
    Since $s < n$, we have $s \leq n-1$, and hence
    \[
        s! \leq (n-1)! \leq \frac{n!}{2} = |A_n|
    \]
    because $n \geq 5$. If $s = 2$, then $K$ is a subgroup of index $2$ in $A_n$, hence normal in $A_n$, contradicting the simplicity of $A_n$. If $s \geq 3$, then Exercise~\ref{ex:3.1} applies to the simple group $A_n$ and rules out a proper subgroup of index $s$.
    This contradiction shows that our assumption $H \neq A_n$ was false. Therefore the only proper subgroup of $S_n$ of index less than $n$ is $A_n$.
\end{solution}

\subsection*{Further Exercises}

Let $X$ be a transitive $G$-set. A \emph{system of imprimitivity} of $X$ is a partition of $X$ which is permuted by the action of $G$. Note that a system of imprimitivity of a $G$-set is itself a $G$-set.

\begin{exercise}
\label{ex:3.9}
Show that there is a bijective correspondence between the set of blocks which contain a given element of $X$ and the set of systems of imprimitivity of $X$. (Observe that when $X$ is finite, this implies that any two elements of $X$ lie in the same number of blocks.)
\end{exercise}
\begin{solution}
    Fix $x \in X$. We construct maps in both directions.
    First let $B$ be a block of $X$ with $x \in B$. Define
    \[
        \mathcal{P}_B = \{gB \mid g \in G\}.
    \]
    We show that $\mathcal{P}_B$ is a system of imprimitivity of $X$.
    The sets in $\mathcal{P}_B$ are either equal or disjoint. Indeed, suppose
    \[
        gB \cap hB \neq \varnothing.
    \]
    Applying $g^{-1}$ gives
    \[
        B \cap g^{-1}hB \neq \varnothing.
    \]
    Since $B$ is a block, this forces $B = g^{-1}hB$. Multiplying by $g$, we get $gB = hB$.
    The sets in $\mathcal{P}_B$ also cover $X$. Let $y \in X$. Since $X$ is transitive, there exists $g \in G$ such that $gx = y$. Since $x \in B$, we have $y = gx \in gB$. Hence every element of $X$ lies in some member of $\mathcal{P}_B$.
    Thus $\mathcal{P}_B$ is a partition of $X$. Moreover, $G$ permutes the pieces, because if $u \in G$, then
    \[
        u(gB) = (ug)B,
    \]
    which is again an element of $\mathcal{P}_B$. Hence $\mathcal{P}_B$ is a system of imprimitivity.
    Conversely, let $\mathcal{P}$ be a system of imprimitivity of $X$. Since $\mathcal{P}$ is a partition, there is a unique piece $B \in \mathcal{P}$ such that $x \in B$. We claim that $B$ is a block. Let $g \in G$. Since $\mathcal{P}$ is permuted by $G$, the set $gB$ is also a piece of $\mathcal{P}$. Since $\mathcal{P}$ is a partition, the two pieces $B$ and $gB$ are either equal or disjoint. Therefore $B$ is a block containing $x$.
    These two constructions undo one another. Starting with a block $B$ containing $x$, the system $\mathcal{P}_B$ contains the piece $B = 1B$. Since $\mathcal{P}_B$ is a partition and $x \in B$, this is the unique piece of $\mathcal{P}_B$ containing $x$. Thus passing from $B$ to $\mathcal{P}_B$ and then taking the piece containing $x$ returns the original block $B$.
    Conversely, start with a system of imprimitivity $\mathcal{P}$, and let $B$ be the unique piece of $\mathcal{P}$ containing $x$. We claim that the system obtained from this block is the original system:
    \[
        \mathcal{P} = \{gB \mid g \in G\}.
    \]
    Since $\mathcal{P}$ is permuted by $G$, every set $gB$ is a piece of $\mathcal{P}$. Conversely, let $P \in \mathcal{P}$. Choose $y \in P$. Since $X$ is transitive, there exists $g \in G$ such that $gx = y$. Now $x \in B$, so $y = gx \in gB$. Also $gB$ is a piece of $\mathcal{P}$. Thus $P$ and $gB$ are two pieces of the partition $\mathcal{P}$ containing $y$, so they must be equal. Hence $P = gB$, and every piece of $\mathcal{P}$ is a translate of $B$.
    Therefore the blocks containing $x$ are in bijective correspondence with the systems of imprimitivity of $X$.
    If $X$ is finite, then for each chosen element $x \in X$, the blocks containing $x$ are in bijection with the same set of systems of imprimitivity. Hence the number of blocks containing $x$ is independent of $x$, so any two elements of $X$ lie in the same number of blocks.
\end{solution}

\begin{exercise}
\label{ex:3.10}
Suppose that the $G$-set $Y$ is an epimorphic image of a $G$-set $X$. Show that there is a $G$-set isomorphism between $Y$ and some system of imprimitivity of $X$.
\end{exercise}
\begin{solution}
    Since $Y$ is an epimorphic image of $X$, there is a surjective $G$-set homomorphism
    \[
        \varphi \colon X \to Y.
    \]
    For each $y \in Y$, let
    \[
        X_y = \varphi^{-1}(y) = \{x \in X \mid \varphi(x) = y\}.
    \]
    Since $\varphi$ is surjective, each $X_y$ is non-empty. The fibers of $\varphi$ form a partition of $X$, so set
    \[
        \mathcal{P} = \{X_y \mid y \in Y\}.
    \]
    We show that $\mathcal{P}$ is a system of imprimitivity. Let $g \in G$ and $y \in Y$. We claim that
    \[
        gX_y = X_{gy}.
    \]
    Indeed, if $x \in X_y$, then $\varphi(x)=y$, and since $\varphi$ is a $G$-set homomorphism,
    \[
        \varphi(gx) = g\varphi(x) = gy.
    \]
    Thus $gx \in X_{gy}$, so $gX_y \subseteq X_{gy}$.
    Conversely, if $z \in X_{gy}$, then $\varphi(z)=gy$. Since $g^{-1}z \in X$, we have
    \[
        \varphi(g^{-1}z) = g^{-1}\varphi(z) = g^{-1}(gy)=y.
    \]
    Hence $g^{-1}z \in X_y$, and therefore $z \in gX_y$. Thus $X_{gy} \subseteq gX_y$, proving $gX_y=X_{gy}$.
    Therefore the action of $G$ permutes the pieces of $\mathcal{P}$, so $\mathcal{P}$ is a system of imprimitivity of $X$.
    Now define
    \[
        \theta \colon Y \to \mathcal{P}, \qquad \theta(y)=X_y.
    \]
    This map is surjective by definition of $\mathcal{P}$. It is injective because distinct elements of $Y$ have disjoint fibers: if $X_y=X_{y'}$, then this common fiber is non-empty, so choose $x$ in it. Then $\varphi(x)=y$ and $\varphi(x)=y'$, hence $y=y'$.
    Finally, $\theta$ is a $G$-set homomorphism. For $g \in G$ and $y \in Y$,
    \[
        \theta(gy)=X_{gy}=gX_y=g\theta(y).
    \]
    Therefore $\theta$ is a $G$-set isomorphism between $Y$ and the system of imprimitivity $\mathcal{P}$.
\end{solution}
If $X$ is a $G$-set, we use $[X]$ to denote the isomorphism class of $X$.

\begin{exercise}
\label{ex:3.11}
Let $G$ be a finite group, and let $S(G)$ be the set of isomorphism classes of finite $G$-sets. Show that we can define sum and product operations on $S(G)$ by $[X] + [Y] = [X \cup Y]$ and $[X][Y] = [X \times Y]$.
\end{exercise}
\begin{solution}
    We must check that these operations do not depend on the chosen representatives of the isomorphism classes.
    Suppose $X \cong X'$ and $Y \cong Y'$ as $G$-sets. Choose $G$-set isomorphisms
    \[
        \alpha \colon X \to X'
        \quad\text{and}\quad
        \beta \colon Y \to Y'.
    \]
    First consider the sum operation. We regard $X \cup Y$ as the disjoint union of $X$ and $Y$, with $G$ acting on each piece in the usual way. Define
    \[
        \alpha \cup \beta \colon X \cup Y \to X' \cup Y'
    \]
    by
    \[
        (\alpha \cup \beta)(z)
        =
        \begin{cases}
            \alpha(z), & z \in X,\\
            \beta(z),  & z \in Y.
        \end{cases}
    \]
    Since $\alpha$ and $\beta$ are bijections, $\alpha \cup \beta$ is a bijection. It also commutes with the $G$-action, because if $z \in X$, then
    \[
        (\alpha \cup \beta)(gz)=\alpha(gz)=g\alpha(z)=g(\alpha \cup \beta)(z),
    \]
    and the same calculation using $\beta$ works when $z \in Y$. Thus $X \cup Y \cong X' \cup Y'$ as $G$-sets, so $[X]+[Y]=[X \cup Y]$ is well-defined.
    Now consider the product operation. The product $X \times Y$ is a $G$-set by the diagonal action
    \[
        g(x,y)=(gx,gy).
    \]
    Define
    \[
        \alpha \times \beta \colon X \times Y \to X' \times Y',
        \qquad
        (\alpha \times \beta)(x,y)=(\alpha(x),\beta(y)).
    \]
    Since $\alpha$ and $\beta$ are bijections, $\alpha \times \beta$ is a bijection. Also,
    \[
        (\alpha \times \beta)(g(x,y))
        =
        (\alpha \times \beta)(gx,gy)
        =
        (\alpha(gx),\beta(gy))
        =
        (g\alpha(x),g\beta(y))
        =
        g(\alpha(x),\beta(y)).
    \]
    Therefore $\alpha \times \beta$ is a $G$-set isomorphism, so $X \times Y \cong X' \times Y'$ as $G$-sets. Hence $[X][Y]=[X \times Y]$ is well-defined.
\end{solution}

\begin{exercise}
\label{ex:3.12}
(cont.) Let $B(G)$ be the set obtained from $S(G)$ by adjoining formal additive inverses of isomorphism classes, in the same way that $\Z$ is obtained from $\N$ by adjoining the additive inverses of positive integers. (The additive identity here will be the isomorphism class of the empty set.) Show that the operations defined above on $S(G)$ extend to give a commutative ring structure on $B(G)$. This ring $B(G)$ is called the \emph{Burnside ring} of $G$.
\end{exercise}
\begin{solution}
    By Exercise~\ref{ex:3.11}, the operations
    \[
        [X]+[Y]=[X\cup Y],
        \qquad
        [X][Y]=[X\times Y]
    \]
    are well-defined on isomorphism classes of finite $G$-sets. The set $B(G)$ is obtained by formally adjoining additive inverses, so every element of $B(G)$ may be represented as a formal difference
    \[
        [X]-[Y],
    \]
    where $X$ and $Y$ are finite $G$-sets.
    Addition is then defined by collecting the positive and negative parts:
    \[
        ([X]-[Y])+([U]-[V])
        =
        [X\cup U]-[Y\cup V].
    \]
    This gives an abelian group under addition. The zero element is $[\varnothing]$, and the additive inverse of $[X]-[Y]$ is $[Y]-[X]$, since
    \[
        ([X]-[Y])+([Y]-[X])
        =
        [X\cup Y]-[Y\cup X]
        =
        [\varnothing]
    \]
    in the formal group obtained from $S(G)$.
    We extend multiplication by the distributive rule:
    \[
        ([X]-[Y])([U]-[V])
        =
        [X\times U]-[X\times V]-[Y\times U]+[Y\times V].
    \]
    This agrees with the original product on $S(G)$ when $Y=V=\varnothing$.
    The usual bijections
    \[
        X\cup Y \cong Y\cup X,
        \qquad
        (X\cup Y)\cup Z \cong X\cup (Y\cup Z)
    \]
    are $G$-set isomorphisms, so addition is commutative and associative. Similarly,
    \[
        X\times Y \cong Y\times X,
        \qquad
        (X\times Y)\times Z \cong X\times (Y\times Z)
    \]
    as $G$-sets, so multiplication is commutative and associative.
    The multiplicative identity is the class $[\{*\}]$ of the one-point $G$-set, where $g*=*$ for every $g\in G$, because
    \[
        X\times \{*\}\cong X
    \]
    as $G$-sets.
    Finally, Cartesian product distributes over disjoint union:
    \[
        X\times (Y\cup Z)
        \cong
        (X\times Y)\cup (X\times Z)
    \]
    as $G$-sets. Hence multiplication distributes over addition in $B(G)$.
    Therefore the operations on $S(G)$ extend to make $B(G)$ a commutative ring.
\end{solution}

\begin{exercise}
\label{ex:3.13}
(cont.) Show that any element of $B(G)$ can be written uniquely as a $\Z$-linear combination of isomorphism classes of transitive $G$-sets.
\end{exercise}
\begin{solution}
    Let $\mathcal{T}$ be the set of isomorphism classes of transitive finite $G$-sets.
    Let $X$ be a finite $G$-set. By Proposition~\ref{prop:3.3}, $X$ has a unique decomposition into orbits:
    \[
        X = O_1\cup \cdots \cup O_r.
    \]
    Each orbit $O_i$ is a transitive $G$-set, and hence
    \[
        [X]=[O_1]+\cdots+[O_r]
    \]
    in $S(G)$. If several of the orbits are isomorphic as $G$-sets, we collect their isomorphism classes together. Thus
    \[
        [X]=n_1[T_1]+\cdots+n_k[T_k],
    \]
    where the $T_i$ are pairwise non-isomorphic transitive $G$-sets and $n_i\in \N$ records the number of orbits of $X$ isomorphic to $T_i$.
    This expression is unique in $S(G)$, because the coefficient $n_i$ is exactly the number of orbits of $X$ whose isomorphism class is $[T_i]$.
    Therefore $S(G)$ is the free commutative monoid on the set $\mathcal{T}$.
    Now let $\alpha\in B(G)$. By construction of $B(G)$, we may write
    \[
        \alpha=[X]-[Y]
    \]
    for finite $G$-sets $X$ and $Y$. Decompose both $X$ and $Y$ into transitive orbits:
    \[
        [X]=n_1[T_1]+\cdots+n_k[T_k],
        \qquad
        [Y]=m_1[T_1]+\cdots+m_k[T_k],
    \]
    after enlarging the list of $T_i$ if necessary. Then
    \[
        \alpha
        =
        [X]-[Y]
        =
        (n_1-m_1)[T_1]+\cdots+(n_k-m_k)[T_k].
    \]
    Hence every element of $B(G)$ is a $\Z$-linear combination of isomorphism classes of transitive $G$-sets.
    Since $S(G)$ is the free commutative monoid on $\mathcal{T}$, adjoining formal additive inverses turns it into the free abelian group on $\mathcal{T}$. Thus an element of $B(G)$ is determined by its integer coefficient attached to each transitive isomorphism class. Therefore no non-trivial $\Z$-linear relation among distinct transitive isomorphism classes is possible, and the expression above is unique.
\end{solution}

\begin{exercise}
\label{ex:3.14}
(cont.) Let $H \leq G$. Show that there is a unique ring homomorphism from $B(G)$ to $\Z$ which sends an isomorphism class $[X]$ to the number of elements of $X$ fixed by $H$.
\end{exercise}
\begin{solution}
    For a finite $G$-set $X$, define
    \[
        X^H=\{x\in X\mid hx=x \text{ for every } h\in H\}.
    \]
    We first check that the rule $[X]\mapsto |X^H|$ is well-defined on isomorphism classes. Suppose that $X\cong Y$ as $G$-sets, and let
    \[
        f\colon X\to Y
    \]
    be a $G$-set isomorphism. If $x\in X^H$, then for every $h\in H$,
    \[
        hf(x)=f(hx)=f(x),
    \]
    so $f(x)\in Y^H$. Applying the same argument to $f^{-1}$ shows that $f$ restricts to a bijection $X^H\to Y^H$. Hence $|X^H|=|Y^H|$.
    Define
    \[
        \varphi_H([X])=|X^H|
    \]
    on $S(G)$. This map respects addition, because
    \[
        (X\cup Y)^H=X^H\cup Y^H
    \]
    as a disjoint union. Hence
    \[
        \varphi_H([X]+[Y])
        =
        |(X\cup Y)^H|
        =
        |X^H|+|Y^H|
        =
        \varphi_H([X])+\varphi_H([Y]).
    \]
    It also respects multiplication, because an element $(x,y)\in X\times Y$ is fixed by every element of $H$ iff both $x$ and $y$ are fixed by every element of $H$. Thus
    \[
        (X\times Y)^H=X^H\times Y^H,
    \]
    and so
    \[
        \varphi_H([X][Y])
        =
        |(X\times Y)^H|
        =
        |X^H|\,|Y^H|
        =
        \varphi_H([X])\varphi_H([Y]).
    \]
    The multiplicative identity $[\{*\}]$ is sent to $1$, since the one point of $\{*\}$ is fixed by every element of $H$.
    Since $B(G)$ is obtained from $S(G)$ by adjoining formal additive inverses, the additive map $\varphi_H$ extends uniquely to $B(G)$ by
    \[
        \varphi_H([X]-[Y])=|X^H|-|Y^H|.
    \]
    The multiplication calculation above shows that this extension is multiplicative, so it is a ring homomorphism $B(G)\to \Z$.
    Finally, this homomorphism is unique. If $\psi\colon B(G)\to \Z$ is any ring homomorphism satisfying $\psi([X])=|X^H|$ for every finite $G$-set $X$, then for every formal difference $[X]-[Y]\in B(G)$,
    \[
        \psi([X]-[Y])
        =
        \psi([X])-\psi([Y])
        =
        |X^H|-|Y^H|.
    \]
    Thus $\psi$ must agree with $\varphi_H$ on every element of $B(G)$, proving uniqueness.
\end{solution}

\begin{exercise}
\label{ex:3.15}
(cont.) Show that any non-zero homomorphism from $B(G)$ to $\Z$ arises from the construction in Exercise~\ref{ex:3.14} and that the intersection of the kernels of all such homomorphisms is zero.
\end{exercise}
\begin{solution}
    By Proposition~\ref{prop:3.4}, every transitive finite $G$-set is isomorphic to $G/K$ for some subgroup $K\leq G$. Here $G/K$ means the set of left cosets of $K$ in $G$, with $G$ acting by left multiplication; it is not being used as a quotient group. Thus, by Exercise~\ref{ex:3.13}, the elements $[G/K]$, with $K$ running over representatives of conjugacy classes of subgroups, generate $B(G)$ additively.\\
    Let $\psi\colon B(G)\to \Z$ be a non-zero ring homomorphism. Since $[G/G]$ is the multiplicative identity of $B(G)$, we have
    \[
        \psi([G/G])^2=\psi([G/G]^2)=\psi([G/G]).
    \]
    The only integers satisfying $n^2=n$ are $0$ and $1$. Hence $\psi([G/G])$ is either $0$ or $1$. If it were $0$, then for every $\alpha\in B(G)$,
    \[
        \psi(\alpha)=\psi([G/G]\alpha)=\psi([G/G])\psi(\alpha)=0,
    \]
    contrary to the assumption that $\psi$ is non-zero. Thus $\psi([G/G])=1$.
    Choose a subgroup $H \leq G$ with $\psi([G/H])\neq 0$ and with $|H|$ minimal among all such subgroups. We claim that $\psi=\varphi_H$, where $\varphi_H$ is the fixed-point homomorphism from Exercise~\ref{ex:3.14}.
    Let $K \leq G$. We decompose the product $G/H\times G/K$ into $G$-orbits using Proposition~\ref{prop:3.3}, where $G$ acts diagonally:
    \[
        u(aH,bK)=(uaH,ubK).
    \]
    Since the first coordinate of an element of $G/H\times G/K$ has the form $aH$, acting by $a^{-1}$ sends $(aH,bK)$ to
    \[
        (H,a^{-1}bK).
    \]
    Hence every $G$-orbit has a representative of the form $(H,gK)$. We now compute the stabilizer of this representative. An element $u\in G$ fixes $(H,gK)$ iff
    \[
        u(H,gK)=(H,gK),
    \]
    which is equivalent to the two conditions
    \[
        uH=H
        \qquad\text{and}\qquad
        ugK=gK.
    \]
    The first condition says $u\in H$. The second condition says $g^{-1}ug\in K$, or equivalently $u\in gKg^{-1}$. Therefore the stabilizer is
    \[
        H\cap gKg^{-1}.
    \]
    By Proposition~\ref{prop:3.4}, the orbit of $(H,gK)$ is therefore isomorphic to
    \[
        G/(H\cap gKg^{-1}).
    \]
    Combining the orbit decomposition from Proposition~\ref{prop:3.3} with the orbit-stabilizer identification from Proposition~\ref{prop:3.4}, we get
    \[
        [G/H][G/K]
        =
        \sum [G/(H\cap gKg^{-1})],
    \]
    where the sum runs over representatives of the $H$-orbits on $G/K$. Applying $\psi$ gives
    \[
        \psi([G/H])\psi([G/K])
        =
        \sum \psi([G/(H\cap gKg^{-1})]).
    \]
    If $H\cap gKg^{-1}$ is a proper subgroup of $H$, then its order is smaller than $|H|$, so its term is zero by the minimal choice of $H$. The only remaining terms occur when
    \[
        H\cap gKg^{-1}=H,
    \]
    equivalently when $H \leq gKg^{-1}$. This is exactly the condition that the coset $gK$ is fixed by every element of $H$. Indeed, $gK$ is fixed by $H$ iff for every $h\in H$,
    \[
        hgK=gK.
    \]
    Multiplying on the left by $g^{-1}$, this is equivalent to
    \[
        g^{-1}hgK=K,
    \]
    which holds iff $g^{-1}hg\in K$. Thus $gK$ is fixed by every element of $H$ iff $g^{-1}Hg\leq K$, equivalently iff $H\leq gKg^{-1}$.
    Hence the number of remaining terms is $|(G/K)^H|$, and the equality above becomes
    \[
        \psi([G/H])\psi([G/K])
        =
        |(G/K)^H|\psi([G/H]).
    \]
    Since $\psi([G/H])\neq 0$, we may cancel it in $\Z$, obtaining
    \[
        \psi([G/K])=|(G/K)^H|=\varphi_H([G/K]).
    \]
    Since the classes $[G/K]$ generate $B(G)$ additively, this proves that $\psi=\varphi_H$. Thus every non-zero homomorphism $B(G)\to \Z$ arises from the construction in Exercise~\ref{ex:3.14}.\\
    It remains to prove that the intersection of the kernels is zero. Let
    \[
        \alpha=\sum_K a_K[G/K]\in B(G),
    \]
    where $K$ runs over one representative from each conjugacy class of subgroups of $G$. Suppose that $\alpha\neq 0$. Choose $K$ with $a_K\neq 0$ and with $|K|$ maximal among all subgroups appearing with non-zero coefficient.
    We evaluate at the fixed-point homomorphism $\varphi_K$. For another subgroup $L$, the set $(G/L)^K$ is non-empty iff some coset $gL$ is fixed by $K$, which is equivalent to
    \[
        g^{-1}Kg \leq L.
    \]
    Thus, if $(G/L)^K$ is non-empty, then $|K|\leq |L|$. If also $a_L\neq 0$, maximality of $|K|$ gives $|L|\leq |K|$, and hence $|L|=|K|$. Therefore $g^{-1}Kg=L$, so $L$ is conjugate to $K$. Since we chose one representative from each conjugacy class, this forces $L=K$.
    Hence all non-zero coefficient terms except the $K$-term make no contribution to $\varphi_K(\alpha)$, and
    \[
        \varphi_K(\alpha)=a_K |(G/K)^K|.
    \]
    The coset $K\in G/K$ is fixed by every element of $K$, so $|(G/K)^K|>0$. Since $a_K\neq 0$, we have $\varphi_K(\alpha)\neq 0$. Therefore $\alpha$ is not in the kernel of $\varphi_K$.
    Thus no non-zero element of $B(G)$ lies in the kernels of all fixed-point homomorphisms, and so the intersection of all these kernels is zero.
\end{solution}

An element $e$ of a ring $R$ is called an \emph{idempotent} if $e^2 = e$.

\begin{exercise}
\label{ex:3.16}
(cont.) Show that if $G$ has no self-normalizing proper subgroup, then $B(G)$ has no idempotents other than the additive and multiplicative identities. (A subgroup $H$ of $G$ is called \emph{self-normalizing} if $N_G(H) = H$. We shall see in Section~11 that there is an important class of groups, namely the nilpotent groups, that have no self-normalizing proper subgroups.)
\end{exercise}
\begin{solution}
    Let $e\in B(G)$ be an idempotent, so $e^2=e$. For every subgroup $H\leq G$, let
    \[
        \varphi_H\colon B(G)\to \Z
    \]
    be the fixed-point homomorphism defined by
    \[
        \varphi_H([X])=|X^H|
    \]
    for every finite $G$-set $X$. Since $\varphi_H$ is a ring homomorphism,
    \[
        \varphi_H(e)^2=\varphi_H(e^2)=\varphi_H(e).
    \]
    Hence $\varphi_H(e)$ is an integer satisfying $n^2=n$, so
    \[
        \varphi_H(e)\in\{0,1\}.
    \]
    We next show that these values do not change when we enlarge a subgroup inside its normalizer. Suppose that
    \[
        H<K\leq N_G(H)
    \]
    and that $|K:H|=p$ is prime. Since $K\leq N_G(H)$, the quotient group $K/H$ acts on $X^H$ for every finite $G$-set $X$. The fixed points of this action are exactly $X^K$. Since $K/H$ has prime order $p$, Corollary~\ref{cor:3.5} shows that each orbit of $K/H$ on $X^H$ has size either $1$ or $p$. The orbits of size $1$ are precisely the points of $X^K$, while every other orbit has size $p$. Hence, for some integer $r\geq 0$,
    \[
        |X^H|=|X^K|+rp.
    \]
    Therefore
    \[
        |X^H|\equiv |X^K|\pmod p.
    \]
    By additivity, the same congruence holds for every element of $B(G)$; in particular,
    \[
        \varphi_H(e)\equiv \varphi_K(e)\pmod p.
    \]
    Since both $\varphi_H(e)$ and $\varphi_K(e)$ are either $0$ or $1$, this congruence forces
    \[
        \varphi_H(e)=\varphi_K(e).
    \]
    Now let $H<G$ be a proper subgroup. By hypothesis, $H$ is not self-normalizing, so
    \[
        H<N_G(H).
    \]
    Thus $N_G(H)/H$ is a non-trivial finite group. Choose a subgroup $K/H$ of prime order in $N_G(H)/H$. Then
    \[
        H<K\leq N_G(H),
    \]
    and the previous paragraph gives $\varphi_H(e)=\varphi_K(e)$. If $K<G$, we repeat the same argument with $K$ in place of $H$. Since $G$ is finite, this process must eventually reach $G$. Hence
    \[
        \varphi_H(e)=\varphi_G(e)
    \]
    for every subgroup $H\leq G$.\\
    Therefore all fixed-point homomorphisms take the same value on $e$, and this common value is either $0$ or $1$. If it is $0$, then $e$ lies in the kernel of every $\varphi_H$, so Exercise~\ref{ex:3.15} gives $e=0$. If it is $1$, then for every $H\leq G$,
    \[
        \varphi_H(e-[G/G])=\varphi_H(e)-\varphi_H([G/G])=1-1=0,
    \]
    since $G/G$ is the one-point $G$-set. Again by Exercise~\ref{ex:3.15}, we get
    \[
        e-[G/G]=0.
    \]
    Thus $e=[G/G]$, the multiplicative identity of $B(G)$. Therefore the only idempotents in $B(G)$ are $0$ and $1$.
\end{solution}

\end{document}
